\documentclass[11pt,a4paper]{article}

% ---------------------------------------------------------------------------
% Paquets
% ---------------------------------------------------------------------------
\usepackage[T1]{fontenc}
\usepackage[utf8]{inputenc}
\usepackage{lmodern}
\usepackage[french]{babel}
\usepackage[a4paper,top=2.5cm,bottom=2.5cm,left=2.5cm,right=2.5cm,headheight=14pt]{geometry}
\usepackage{graphicx}
\usepackage{booktabs}
\usepackage{tabularx}
\usepackage{longtable}
\usepackage{array}
\usepackage{float}
\usepackage{caption}
\usepackage{xcolor}
\usepackage{listings}
\usepackage{fancyhdr}
\usepackage{enumitem}
\usepackage{textcomp}
\usepackage[colorlinks=true,linkcolor={blue!50!black},urlcolor={blue!50!black}]{hyperref}

\graphicspath{{figures/}{logos/}}

% En-tete / pied de page (comme le compte rendu du TP1)
\pagestyle{fancy}
\fancyhf{}
\fancyhead[L]{TP2 CI/CD sécurisé}
\fancyhead[R]{Compte rendu}
\fancyfoot[C]{\thepage}
\renewcommand{\headrulewidth}{0.4pt}

% Listings
\lstset{
  basicstyle=\ttfamily\footnotesize,
  frame=single,
  breaklines=true,
  breakatwhitespace=false,
  columns=fixed,
  keepspaces=true,
  showstringspaces=false,
  upquote=true,
  tabsize=4,
  xleftmargin=0pt,
  framexleftmargin=0pt,
  aboveskip=6pt,
  belowskip=8pt,
  literate={é}{{\'e}}1 {è}{{\`e}}1 {à}{{\`a}}1 {ê}{{\^e}}1 {ô}{{\^o}}1 {î}{{\^i}}1 {û}{{\^u}}1 {ç}{{\c c}}1 {É}{{\'E}}1 {→}{{$\rightarrow$}}1
}
\renewcommand{\lstlistingname}{Listing}
\captionsetup[lstlisting]{labelsep=endash}

\newcolumntype{L}{>{\raggedright\arraybackslash}X}
\setlist[itemize]{label=---,leftmargin=*}

\newcommand{\code}[1]{\texttt{#1}}
\setlength{\emergencystretch}{3em}

% Figure encadree : \figshot{fichier}{largeur}{legende}{etiquette}
\newcommand{\figshot}[4]{%
\begin{figure}[H]
\centering
\fbox{\includegraphics[width=#2\textwidth]{#1}}
\caption{#3}
\label{#4}
\end{figure}}

% ===========================================================================
\begin{document}

% ---------------------------------------------------------------------------
% Page de garde
% ---------------------------------------------------------------------------
\begin{titlepage}
\thispagestyle{empty}
\noindent
\begin{minipage}[c]{0.2\textwidth}
  \includegraphics[width=\linewidth]{insat.jpg}
\end{minipage}%
\hfill
\begin{minipage}[c]{0.56\textwidth}
  \centering\small
  RÉPUBLIQUE TUNISIENNE\\
  Ministère de l'Enseignement Supérieur\\
  et de la Recherche Scientifique\\
  Université de Carthage\\
  \textbf{Institut National des Sciences Appliquées}\\
  \textbf{et de Technologie}
\end{minipage}%
\hfill
\begin{minipage}[c]{0.2\textwidth}
  \includegraphics[width=\linewidth]{ucar.jpg}
\end{minipage}

\vspace{2.6cm}
\begin{center}
  {\large\textsc{DevSecOps}}\\[0.4cm]
  Spécialité : Génie Logiciel\\
  Niveau : 5\textsuperscript{e} année

  \vspace{1.4cm}
  \rule{\textwidth}{0.4pt}\\[0.5cm]
  {\huge\bfseries TP2 -- CI/CD sécurisé}\\[0.4cm]
  {\Large\bfseries Intégration de la Sécurité dans un Pipeline CI/CD}\\[0.3cm]
  \rule{\textwidth}{0.4pt}

  \vspace{1.4cm}
  {\large Sécurisation d'une application web Flask conteneurisée}\\[0.1cm]
  Bandit, pip-audit, OWASP Dependency-Check, OWASP ZAP et GitHub Actions

  \vspace{1.6cm}
  Présenté par :\\[0.15cm]
  {\large\bfseries Yassine Dahmoul}\\[0.05cm]
  {\large\bfseries Mohamed Sahbi Ben Kalaia}\\[0.05cm]
  {\large\bfseries Mohamed Dhia Medini}\\[0.05cm]
  {\large\bfseries Takoua Ayadi}


  \vfill
  \rule{0.55\textwidth}{0.4pt}\\[0.3cm]
  \textbf{Année Universitaire : 2026 -- 2027}
\end{center}
\end{titlepage}

\setcounter{page}{1}
\tableofcontents
\newpage

% ---------------------------------------------------------------------------
\section{Contexte et Objectifs}
% ---------------------------------------------------------------------------
Au TP1, les scans de sécurité étaient lancés \emph{à la main}, une seule fois. Rien n'oblige
quelqu'un à les relancer après chaque modification, chaque nouveau commit peut introduire une
faille, et de nouvelles vulnérabilités sont publiées chaque jour sur des bibliothèques qui n'ont
pas changé. Ce TP automatise ces contrôles dans un \textbf{pipeline CI/CD} qui joue le rôle de
\textbf{garde-fou} (\emph{guardrail}) : à chaque modification, les mêmes contrôles s'exécutent
automatiquement, et leur résultat \textbf{décide} de ce qui peut être fusionné et publié.

Trois familles de contrôles, complémentaires, sont intégrées au pipeline :

\begin{table}[H]
\centering
\begin{tabularx}{\textwidth}{@{}lLL@{}}
\toprule
\textbf{Approche} & \textbf{Ce qui est analysé} & \textbf{Outil utilisé} \\
\midrule
\textbf{SAST} (\emph{Static Application Security Testing}) & Notre code source, sans l'exécuter & Bandit \\
\textbf{SCA} (\emph{Software Composition Analysis}) & Les bibliothèques tierces (directes et transitives) & pip-audit, OWASP Dependency-Check \\
\textbf{DAST} (\emph{Dynamic Application Security Testing}) & L'application \emph{en cours d'exécution}, attaquée de l'extérieur & OWASP ZAP \\
\bottomrule
\end{tabularx}
\end{table}

\noindent Objectifs principaux (sujet du TP) :
\begin{itemize}
  \item Automatiser les tests de sécurité dans un processus de développement continu.
  \item Comprendre le rôle d'un pipeline CI/CD comme \textbf{garde-fou de sécurité}.
  \item Configurer des \textbf{règles de validation} qui bloquent les déploiements non sécurisés.
\end{itemize}

Le travail suit les quatre activités du sujet :

\begin{table}[H]
\centering
\begin{tabularx}{\textwidth}{@{}lL@{}}
\toprule
\textbf{Activité} & \textbf{Réalisation} \\
\midrule
1. Pipeline de base & Compilation, tests unitaires et construction de l'image Docker (Partie~I) \\
2. SAST et SCA & Bandit, pip-audit et Dependency-Check avec des seuils bloquants (Partie~II) \\
3. DAST & Scan OWASP ZAP authentifié sur un environnement éphémère (Partie~III) \\
   & Porte de sécurité, protection de la branche et publication conditionnelle (Partie~IV) \\
4. Validation & Une pull request volontairement vulnérable (pipeline en échec), puis une pull request corrigée (pipeline vert) (Partie~V) \\
\bottomrule
\end{tabularx}
\end{table}

% ---------------------------------------------------------------------------
\section{Environnement Technique}
% ---------------------------------------------------------------------------
\begin{table}[H]
\centering
\begin{tabularx}{\textwidth}{@{}lL@{}}
\toprule
\textbf{Élément} & \textbf{Détail} \\
\midrule
Plateforme CI/CD & GitHub Actions (exécuteurs \code{ubuntu-latest}) \\
Dépôt & \href{https://github.com/YassineDHML/tp2-devsecops-cicd-pipeline}{\code{YassineDHML/tp2-devsecops-cicd-pipeline}} \\
Application & \emph{Notes App} (Python 3.13, Flask), version corrigée du TP1 \\
Tests & pytest 9.2.0, pytest-cov 7.1.0 \\
SAST & Bandit 1.9.4 \\
SCA & pip-audit 2.10.1 (base PyPA / OSV) ; OWASP Dependency-Check 13.0.0 (base NVD, image Docker) \\
DAST & OWASP ZAP 2.17.0 (image Docker \code{ghcr.io/zaproxy/zaproxy:2.17.0}), \emph{Automation Framework} \\
Conteneurisation & Docker : image de l'application, réseau et conteneurs éphémères du DAST \\
Registre d'images & GitHub Container Registry (\code{ghcr.io}) \\
\bottomrule
\end{tabularx}
\end{table}

\noindent Les images des scanners sont épinglées par version, et toutes les actions GitHub par
\textbf{empreinte de commit} : un tag peut être déplacé par un attaquant qui prend le contrôle d'une
action, pas une empreinte.

% ---------------------------------------------------------------------------
\section{Partie I -- Pipeline de base : compiler et tester}
% ---------------------------------------------------------------------------
\subsection{Structure du projet}

\begin{lstlisting}
tp2-devsecops-cicd-pipeline/
 .github/workflows/ci.yml     # le pipeline (un seul fichier)
 app/                         # application Flask issue du TP1
 tests/                       # tests de l'application et du script de gate
 scripts/
   dast.sh                    # environnement ephemere, scan ZAP, gate, destruction
   security_report.py         # resumes des rapports et regle de blocage du DAST
 zap/automation.yaml          # plan ZAP : scan authentifie
 reports/                     # preuves d'execution (PR vulnerable / PR corrigee)
 Dockerfile  requirements.txt  requirements-dev.txt
 RAPPORT.md                   # rapport complet au format Markdown
\end{lstlisting}

Le dépôt est public sur GitHub (Figure~\ref{fig:repo}) : dossiers du projet, dernier commit validé par
le pipeline et lien vers la pull request ouverte.

\figshot{t_gh_repo.png}{0.92}{Dépôt GitHub (branche de la pull request corrigée) : fichiers du projet et dernier commit validé.}{fig:repo}

\subsection{Application et tests}
L'application est la version corrigée du TP1 (\emph{Notes App} : comptes utilisateurs, notes,
recherche, import YAML, préférences, deux outils réseau). Pour que le pipeline ait quelque chose à
vérifier, une suite \code{pytest} a été ajoutée (client de test Flask, base SQLite temporaire pour
chaque test) :

\begin{itemize}
  \item \textbf{comportement fonctionnel} : inscription, connexion, déconnexion, notes, recherche,
        import, préférences ;
  \item \textbf{propriétés de sécurité corrigées au TP1} : requêtes paramétrées, échappement HTML,
        mots de passe hachés, cookie \code{HttpOnly} / \code{SameSite}, YAML sûr, ping sans shell,
        protection anti-SSRF, en-têtes de sécurité. Si l'une de ces corrections était annulée par
        erreur, un test échouerait ;
  \item \textbf{script de gate} : \code{tests/test\_security\_report.py} vérifie la règle de blocage
        elle-même, car un garde-fou défaillant est pire que pas de garde-fou.
\end{itemize}

Au total, \textbf{79 tests} (dont 14 pour le script de gate) et \textbf{99~\%} de couverture
(lignes et branches) de \code{app/app.py}. Le pipeline exige 90~\% (Figure~\ref{fig:pytest}).

\figshot{t_log_pytest.png}{0.85}{Étape \og Tests unitaires et couverture \fg{} du job \emph{Build et tests unitaires} : 79 tests réussis, couverture de 98,96~\%.}{fig:pytest}

\subsection{Étude du fichier \texttt{.github/workflows/ci.yml}}

\begin{table}[H]
\centering
\small
\begin{tabularx}{\textwidth}{@{}lL@{}}
\toprule
\textbf{Bloc du fichier} & \textbf{Rôle} \\
\midrule
\code{on:} \code{pull\_request}, \code{push}, \code{schedule}, \code{workflow\_dispatch} & Contrôle chaque PR \textbf{avant} la fusion (c'est là que le garde-fou agit), le code fusionné sur \code{main}, et une fois par semaine pour détecter les nouvelles CVE sans nouveau commit \\
\code{permissions: contents: read} & Moindre privilège : le jeton ne peut que lire le dépôt. Seul le job \code{release} obtient \code{packages: write} \\
\code{concurrency} & Une nouvelle poussée sur une PR annule l'exécution précédente. Jamais sur \code{main} : une publication en cours n'est pas interrompue \\
\code{env} & Seuils de blocage (\code{BANDIT\_FAIL\_LEVEL}, \code{DC\_FAIL\_ON\_CVSS}) définis une seule fois \\
\code{uses: actions/...@<empreinte>} & Actions épinglées par empreinte de commit (protection de la chaîne d'approvisionnement) \\
\code{timeout-minutes} & Un job bloqué ne consomme pas de minutes indéfiniment \\
\code{if: always()} & Exécute l'étape même après un échec : publication des rapports, destruction de l'environnement éphémère \\
\code{needs} & Ordonne les jobs : le DAST attend l'image, la release attend la porte de sécurité \\
\bottomrule
\end{tabularx}
\caption{Lecture commentée du fichier de configuration du pipeline (fichier complet en annexe)}
\end{table}

\subsection{Les deux jobs de la phase 1}
\paragraph{\code{build-test}.} Python n'est pas compilé, mais deux étapes jouent ce rôle :
\code{python -m compileall} compile en byte-code \textbf{tous} les fichiers, y compris ceux
qu'aucun test n'importe, et \code{pytest -{}-cov -{}-cov-fail-under=90} exécute les tests et échoue si
la couverture est insuffisante. Les rapports JUnit et de couverture sont publiés en artefacts.

\paragraph{\code{build-image}.} \code{docker build} construit \textbf{le livrable}, ce qui vérifie que
le \code{Dockerfile} et \code{requirements.txt} sont cohérents. L'image est étiquetée avec le SHA du
commit, exportée (\code{docker save}) et publiée comme artefact d'une journée. Le DAST la teste et la
release la publie : \textbf{c'est exactement la même image}, pas une reconstruction.

\subsection{Construction progressive du pipeline}
L'historique de \code{main} suit les activités du sujet. Chaque commit qui contient un pipeline est
validé par ses propres contrôles : 2 contrôles pour le pipeline de base, 5 avec le SAST et la SCA, 8 avec
le DAST, la porte de sécurité et la publication (Figure~\ref{fig:commits}).

\figshot{t_gh_commits.png}{0.92}{Historique de \code{main} : le nombre de contrôles réussis par commit passe de 2 à 5 puis à 8.}{fig:commits}

% ---------------------------------------------------------------------------
\section{Partie II -- SAST et SCA avec seuils bloquants}
% ---------------------------------------------------------------------------
\subsection{Politique de blocage}
Le sujet demande de bloquer le pipeline en cas de vulnérabilité \textbf{critique ou de haute
sévérité}. Chaque outil a sa propre échelle ; le tableau montre comment la règle est traduite.

\begin{table}[H]
\centering
\small
\begin{tabularx}{\textwidth}{@{}lLLL@{}}
\toprule
\textbf{Outil} & \textbf{Échelle native} & \textbf{Seuil de blocage} & \textbf{Mécanisme} \\
\midrule
\textbf{Bandit} (SAST) & Low, Medium, High (et confiance) & \textbf{High} & \code{bandit -{}-severity-level high} : code de sortie 1 \\
\textbf{pip-audit} (SCA) & Pas de score CVSS dans la base PyPA / OSV & \textbf{Toute vulnérabilité connue} & Code de sortie 1 \\
\textbf{Dependency-Check} (SCA) & CVSS : Medium $<$ 7 $\leq$ High $<$ 9 $\leq$ Critical & \textbf{CVSS $\geq$ 7} & \code{-{}-failOnCVSS 7} : code de sortie 15 \\
\textbf{OWASP ZAP} (DAST) & Informational, Low, Medium, High (pas de \emph{Critical}) & \textbf{High}, le niveau le plus élevé & \code{security\_report.py zap -{}-fail-on high} \\
\bottomrule
\end{tabularx}
\caption{Traduction de la règle \og critique ou haute sévérité \fg{} pour chaque outil}
\end{table}

\noindent Les constatations de niveau inférieur ne bloquent pas, mais elles sont \textbf{affichées}
dans l'onglet \emph{Summary} du run et conservées dans les rapports. Pour pip-audit, sans score CVSS,
la règle est plus stricte que celle du sujet : toute vulnérabilité connue bloque, ce qui est peu
pénalisant car l'outil indique toujours la version corrigée.

\subsection{SAST -- Bandit}
Le travail est fait en deux temps : d'abord les rapports \textbf{complets}, tous niveaux confondus
(\code{-{}-exit-zero} : cette étape ne bloque jamais), puis le \textbf{gate}, qui relance Bandit avec le
seuil et échoue si une constatation l'atteint. Les rapports sont ainsi toujours produits, et le
journal du gate affiche directement le code fautif.

\begin{lstlisting}[caption={Étapes SAST du pipeline}]
- name: Rapports complets (tous niveaux, non bloquants)
  run: |
    bandit -r app -f json -o ci-reports/bandit.json --exit-zero
    bandit -r app -f html -o ci-reports/bandit.html --exit-zero
    python scripts/security_report.py bandit ci-reports/bandit.json
- name: "GATE SAST : echec si une vulnerabilite de severite HIGH est detectee"
  run: bandit -r app --severity-level "$BANDIT_FAIL_LEVEL"
\end{lstlisting}

Sur l'application saine, Bandit signale 2 constatations \emph{Low} (\code{B404}, import de
\code{subprocess}, et \code{B603}, appel sans shell), analysées et acceptées au TP1 : le gate passe.

\subsection{SCA -- pip-audit}
pip-audit interroge la base \textbf{PyPA / OSV} (avis GitHub, PyPI) sur les versions figées de
\code{requirements.txt}, dépendances transitives comprises, et indique la \textbf{version corrigée}.

\begin{lstlisting}[caption={Étapes pip-audit du pipeline}]
- name: Rapport JSON et resume (non bloquants)
  run: |
    pip-audit -r requirements.txt --no-deps --disable-pip -f json -o ci-reports/pip-audit.json || true
    python scripts/security_report.py pip-audit ci-reports/pip-audit.json
- name: "GATE SCA (1/2) : echec si une dependance a une vulnerabilite connue"
  run: pip-audit -r requirements.txt --no-deps --disable-pip --desc off
\end{lstlisting}

\subsection{SCA -- OWASP Dependency-Check}
\begin{lstlisting}[caption={Lancement d'OWASP Dependency-Check dans le pipeline}]
docker run --rm \
  -v "$PWD:/src:ro" \
  -v "$HOME/dc-data:/usr/share/dependency-check/data" \
  -v "$PWD/ci-reports/dependency-check:/report" \
  owasp/dependency-check:13.0.0 \
  --project "TP2 Notes App" --scan /src \
  --exclude "**/.git/**" --exclude "**/ci-reports/**" \
  --enableExperimental --disableOssIndex \
  --nvdDatafeed "https://dependency-check.github.io/DependencyCheck_Builder/nvd_cache/" \
  --format HTML --format JSON --out /report \
  --failOnCVSS "$DC_FAIL_ON_CVSS"
\end{lstlisting}

\begin{itemize}
  \item \code{-{}-failOnCVSS 7} : \textbf{règle de blocage}, code de sortie 15 dès qu'une CVE a un
        score supérieur ou égal à 7.
  \item \code{-{}-enableExperimental} : l'analyseur de \code{requirements.txt} est expérimental ; sans
        cette option, aucune dépendance Python n'est détectée (constaté au TP1).
  \item \code{-{}-nvdDatafeed} : miroir officiel de la base NVD (le téléchargement par l'API est très
        lent sans clé). La base est conservée d'un run à l'autre avec \code{actions/cache}, sous une
        clé \textbf{hebdomadaire}.
  \item Le code est monté en \textbf{lecture seule} (\code{:ro}) et l'image est épinglée en version
        \code{13.0.0}.
\end{itemize}

Sur l'application saine : 14 dépendances identifiées dans \code{requirements.txt}, \textbf{0 CVE}.

\subsection{Lisibilité des résultats}
Le script \code{scripts/security\_report.py} lit le rapport JSON de chaque outil et en écrit un résumé
(comptes par niveau, détail des constatations) dans l'onglet \emph{Summary} du run. Chaque job publie
aussi ses rapports complets (JSON et HTML) comme \textbf{artefacts}. Pour le DAST, c'est ce script qui
applique la règle de blocage ; il est donc testé, y compris le cas \og rapport absent ou vide \fg{} :
le gate \textbf{échoue fermé}, car sans rapport exploitable la vérification ne peut pas être validée.

\subsection{Un constat réel : une dépendance du TP1 devenue vulnérable}\label{sec:certifi}
En préparant ce TP, pip-audit a été lancé \emph{en local} sur le \code{requirements.txt} du TP1, sans
aucune modification du code (Figure~\ref{fig:certifi}).

\figshot{t_log_certifi.png}{0.85}{pip-audit sur le \code{requirements.txt} du TP1 : \code{certifi} 2026.7.22 est devenu vulnérable.}{fig:certifi}

L'avis (GHSA-6q48-6c7w-9gw2) a été publié le \textbf{20 octobre 2026}, après le TP1 : les versions de
\code{certifi} antérieures à 2026.10.19 contiennent des certificats racine \emph{emSign} dont des
certificats intermédiaires ont des clés privées compromises. La dépendance a été mise à jour
(\code{certifi==2026.10.23}). Deux enseignements :

\begin{enumerate}[leftmargin=*]
  \item \textbf{Le code n'a pas changé, mais le risque, si.} C'est la raison d'être de l'exécution
        planifiée du pipeline : un projet \og terminé \fg{} doit continuer à être contrôlé.
  \item \textbf{Les outils SCA ne voient pas tous la même chose.} Dependency-Check, lancé sur la même
        version, a rapporté \textbf{0 vulnérabilité} : son jeu de données (NVD) ne référence que les
        CVE, or cet avis n'avait pas encore d'identifiant CVE. pip-audit s'appuie sur la base
        PyPA / OSV, qui reçoit les avis GitHub avant l'attribution d'une CVE. D'où l'intérêt de
        \textbf{combiner deux sources}.
\end{enumerate}

% ---------------------------------------------------------------------------
\section{Partie III -- DAST sur un environnement éphémère}
% ---------------------------------------------------------------------------
\subsection{Pourquoi un DAST}
Le SAST lit le code ; le DAST \textbf{attaque l'application qui tourne}, comme le ferait un
attaquant, sans connaître le code. Il voit ce que le SAST ne peut pas voir : la configuration réelle
(en-têtes, cookies), l'effet d'une entrée sur la page produite, ce qu'une route répond sans
authentification. Le prix : il est plus lent et il faut une application déployée.

\subsection{L'environnement éphémère}
\og Éphémère \fg{} signifie créé pour un seul scan, \textbf{identique à chaque fois}, et détruit
ensuite. Le script \code{scripts/dast.sh} en pilote le cycle de vie ; le workflow appelle chaque étape
séparément pour qu'elles soient visibles dans l'interface (Figure~\ref{fig:jobdast}).

\begin{table}[H]
\centering
\small
\begin{tabularx}{\textwidth}{@{}llL@{}}
\toprule
\textbf{Étape du job \code{dast}} & \textbf{Commande} & \textbf{Ce qui se passe} \\
\midrule
Charger l'image & \code{docker load} & Récupère l'image construite par \code{build-image} : c'est le livrable, pas une reconstruction \\
\textbf{Déployer} & \code{dast.sh deploy} & Crée un réseau Docker dédié, lance le conteneur applicatif, attend qu'il réponde, crée un compte de test \\
\textbf{Scanner} & \code{dast.sh scan} & Lance ZAP dans un second conteneur du même réseau (cible \code{http://app:5000}) \\
\textbf{Gate} & \code{dast.sh gate} & Lit le rapport ZAP, échoue s'il contient une alerte High \\
\textbf{Détruire} & \code{dast.sh destroy} & Supprime conteneur et réseau, même si le scan a échoué (\code{if: always()}) \\
\bottomrule
\end{tabularx}
\caption{Cycle de vie de l'environnement éphémère}
\end{table}

\figshot{t_gh_job_dast.png}{0.9}{Job \emph{DAST -- OWASP ZAP} : chargement de l'image, déploiement, scan, gate et destruction de l'environnement.}{fig:jobdast}

Propriétés de sécurité de cet environnement :

\begin{table}[H]
\centering
\small
\begin{tabularx}{\textwidth}{@{}lL@{}}
\toprule
\textbf{Propriété} & \textbf{Mise en œuvre} \\
\midrule
Isolation & Réseau Docker dédié ; le port de l'application n'est publié que sur \code{127.0.0.1} du runner \\
Secrets jetables & \code{SECRET\_KEY} et mot de passe du compte de test générés avec \code{openssl rand -hex} à chaque exécution, stockés dans un fichier en mode \code{0600} supprimé à la fin ; \textbf{rien dans le dépôt ni dans l'image} \\
État vierge & Base SQLite créée dans le conteneur et détruite avec lui \\
Pas de compte par défaut & \code{ADMIN\_PASSWORD} n'est pas défini \\
Processus non privilégié & Le conteneur s'exécute avec l'utilisateur \code{appuser} (durcissement du TP1) \\
Machine jetable & Le job tourne sur un exécuteur GitHub neuf, détruit à la fin \\
\bottomrule
\end{tabularx}
\end{table}

\subsection{Le plan ZAP (Automation Framework)}
Un fichier YAML décrit le scan, ce qui le rend versionné, relisible et reproductible
(\code{zap/automation.yaml}, en annexe).

\begin{table}[H]
\centering
\small
\begin{tabularx}{\textwidth}{@{}lL@{}}
\toprule
\textbf{Élément du plan} & \textbf{Rôle} \\
\midrule
\code{contexts} & Périmètre du scan. \code{/logout} est \textbf{exclu} : sinon le spider suivrait le lien \og Déconnexion \fg{} et fermerait sa propre session \\
\code{authentication: form} & ZAP se connecte via \code{/login} avec le compte jetable (variables d'environnement \code{ZAP\_USER} et \code{ZAP\_PASSWORD}) ; la session est reconnue par la présence du lien \og Deconnexion \fg \\
\code{alertFilter} & Requalifie en \og faux positif \fg{} \textbf{une} alerte précise, après triage (voir \S\ref{sec:fp}) \\
\code{requestor} & Déclare \code{/tools/ping} et \code{/tools/fetch}, que \textbf{aucune page HTML ne référence} : sans cela, le spider ne les découvrirait jamais \\
\code{spider} + test \code{stats} & Explore l'application \textbf{authentifiée}. \textbf{Garde-fou de couverture} : le job échoue si moins de 8 URL sont découvertes (13 à 14 observées), pour qu'un scan vide ne soit jamais validé \\
\code{activeScan} & \textbf{Attaque} : injections SQL et de commande, XSS, traversée de chemin\ldots{} Politique par défaut, durée maximale de 12 minutes \\
\code{report} (2 fois) & Rapport JSON (lu par le gate) et rapport HTML (lu par un humain) \\
\bottomrule
\end{tabularx}
\caption{Les éléments du plan ZAP}
\end{table}

\noindent Un scan \textbf{authentifié} est important : presque toute l'application est derrière la
connexion. La Figure~\ref{fig:zapplan} montre le journal du plan : les deux routes déclarées, les 14 URL
découvertes, le garde-fou de couverture satisfait, puis la fin du scan actif.

\figshot{t_log_zap_plan.png}{0.9}{Journal du plan ZAP (PR corrigée) : spider, garde-fou de couverture, scan actif.}{fig:zapplan}

\subsection{Règle de blocage}
ZAP classe ses alertes en \emph{Informational}, \emph{Low}, \emph{Medium} et \emph{High} ; il n'a pas
de niveau \emph{Critical}. \textbf{High est donc le niveau \og critique \fg{}} du DAST, et c'est celui
qui bloque. ZAP n'ayant pas d'option équivalente à \code{-{}-failOnCVSS}, la règle est appliquée par
\code{security\_report.py zap -{}-fail-on high} sur le rapport JSON. Le script échoue \textbf{fermé} :
rapport absent, illisible ou sans site analysé, il renvoie le code 2 et le job échoue.

\subsection{Résultats sur l'application saine}
Le scan de l'application saine donne \textbf{0 alerte High}, 4 Medium et 3 Informational
(Figure~\ref{fig:zappr2}) : le gate passe. Les alertes Medium ne bloquent pas mais sont de vraies
observations, à analyser (triage) :

\figshot{r_zap_pr2.png}{0.88}{Rapport ZAP de l'application corrigée : aucune alerte High.}{fig:zappr2}

\begin{table}[H]
\centering
\small
\begin{tabularx}{\textwidth}{@{}LlL@{}}
\toprule
\textbf{Alerte (niveau)} & \textbf{CWE} & \textbf{Analyse} \\
\midrule
Absence of Anti-CSRF Tokens (Medium) & 352 & \textbf{Vraie faiblesse} : pas de jeton CSRF dans les formulaires. Le cookie \code{SameSite=Lax} limite le risque. Correctif : Flask-WTF \\
CSP: Failure to Define Directive with No Fallback (Medium) & 693 & \code{frame-ancestors} et \code{form-action} n'héritent pas de \code{default-src}. Correctif : les ajouter à la politique \\
CSP: style-src unsafe-inline (Medium) & 693 & Le CSS est dans une balise \code{<style>}. Correctif : fichier statique, retrait de \code{'unsafe-inline'} \\
HTTP Only Site (Medium) & 311 & \textbf{Attendu} : l'environnement éphémère sert l'application en HTTP. En production : HTTPS, HSTS, cookie \code{Secure} \\
3 alertes Informational & -- & Informations de ZAP sur son propre fonctionnement \\
\bottomrule
\end{tabularx}
\caption{Triage des alertes de niveau Medium (hors périmètre du TP, laissées en l'état)}
\end{table}

\subsection{Un faux positif rencontré : triage et filtre ciblé}\label{sec:fp}
Pendant la validation, un run de la PR corrigée a été \textbf{bloqué par le DAST}, alors que le code
applicatif n'avait pas changé depuis le run précédent, entièrement vert : le commit ne modifiait que la
documentation (Figures~\ref{fig:runfp} et~\ref{fig:logfp}).

\figshot{t_gh_run_fp.png}{0.92}{Exécution n\textsuperscript{o}10 de la PR corrigée : seul le job DAST (et donc la porte de sécurité) est en échec.}{fig:runfp}

\figshot{t_log_fp.png}{0.85}{Le gate DAST signale une injection SQL (High) sur \code{POST /preferences}, paramètre \code{theme}.}{fig:logfp}

Le gate a fait son travail (une alerte High bloque). Reste à décider si l'alerte est fondée : c'est le
\textbf{triage}, en trois temps.

\begin{enumerate}[leftmargin=*]
  \item \textbf{Lire l'alerte} (Figure~\ref{fig:zapfp}). ZAP a utilisé sa règle de détection
        \emph{booléenne} : il compare les pages obtenues avec \code{dark' AND '1'='1' -- } et
        \code{dark' AND '1'='2' -- }. Il n'y a ni message d'erreur SQL ni donnée extraite : la seule
        preuve est une différence entre deux pages, avec une confiance \emph{Medium}.
  \item \textbf{Lire le code.} Dans \code{/preferences}, \code{theme} est comparé à la liste blanche
        \code{\{"light", "dark"\}} puis rangé dans le cookie de session. \textbf{Aucune requête SQL} :
        la route n'utilise même pas la base de données.
  \item \textbf{Expliquer le mécanisme} (hypothèse la plus probable). La page affichée dépend de la
        valeur stockée dans la \textbf{session}, que les requêtes du scanner modifient en permanence :
        le scan actif utilise plusieurs fils d'exécution qui partagent la même session. Deux réponses
        comparées par ZAP peuvent donc différer parce que l'état a changé entre-temps, pas parce que la
        saisie a modifié une requête SQL. Sans requête SQL, il n'y a de toute façon pas d'injection SQL
        possible.
\end{enumerate}

\figshot{r_zap_fp.png}{0.92}{Détail de l'alerte dans le rapport ZAP : seule une différence de pages, aucune preuve d'exécution SQL.}{fig:zapfp}

\noindent\textbf{Conclusion : faux positif.} Le traitement est un \textbf{filtre étroit, documenté et
gardé par un test} :

\begin{itemize}
  \item un job \code{alertFilter} de \code{zap/automation.yaml} requalifie cette alerte en
        \emph{False Positive} pour \textbf{une seule combinaison} : règle 40018, URL, paramètre
        \code{theme}, méthode POST. La règle reste active partout ailleurs. Un essai préalable sur une
        autre alerte a confirmé que le filtre n'écarte que l'instance visée (5 instances ramenées à 4) ;
  \item le test \code{test\_preferences\_never\_touch\_the\_database} \textbf{verrouille l'hypothèse} :
        si une évolution du code faisait utiliser SQL à cette route, le test échouerait et le filtre
        devrait être retiré ;
  \item les preuves du run bloqué sont conservées dans \code{reports/pr2-faux-positif/}.
\end{itemize}

\noindent Ce qu'on ne fait \textbf{pas} : relancer le job jusqu'à ce qu'il passe (cela masquerait aussi
une vraie vulnérabilité intermittente), baisser le seuil global ou désactiver la règle d'injection SQL.
Le run suivant, avec le filtre, est repassé au vert (Figure~\ref{fig:actions}).

\figshot{t_gh_actions.png}{0.9}{Liste des exécutions du pipeline : \#5 en échec (PR vulnérable), \#6 vert, \#7 à \#9 annulées par un \emph{push} plus récent, \#10 bloquée par le faux positif, \#11 verte.}{fig:actions}

% ---------------------------------------------------------------------------
\section{Partie IV -- Porte de sécurité, protection de la branche et publication}
% ---------------------------------------------------------------------------
\subsection{Le job \texttt{security-gate}}
\begin{lstlisting}[caption={Porte de sécurité : agrégat de tous les contrôles}]
security-gate:
  name: Security gate
  needs: [build-test, build-image, sast, sca-pip-audit, sca-dependency-check, dast]
  if: always()
  steps:
    - name: Bilan des controles
      env:
        NEEDS: ${{ toJSON(needs) }}
      run: |
        ...
        if echo "$NEEDS" | jq -e 'to_entries | any(.value.result != "success")' > /dev/null; then
          echo "::error::Au moins un controle n'est pas vert : version non securisee, publication bloquee."
          exit 1
        fi
\end{lstlisting}

Deux détails importants :
\begin{itemize}
  \item \textbf{\code{if: always()}}. Par défaut, un job dont une dépendance a échoué est
        \emph{ignoré}. Or GitHub considère un contrôle obligatoire ignoré comme \textbf{satisfait} :
        sans \code{if: always()}, un échec du SAST laisserait passer la fusion. Le job doit donc
        s'exécuter dans tous les cas et échouer lui-même.
  \item \textbf{\code{result != "success"}}. Le gate échoue non seulement si un contrôle est en échec,
        mais aussi s'il est \emph{ignoré} ou \emph{annulé} : seul un contrôle explicitement vert compte.
\end{itemize}

\subsection{Protection de la branche \texttt{main}}
Un pipeline qui échoue ne protège rien si l'on peut fusionner quand même. La branche \code{main} est
donc protégée (règle créée avec l'API REST de GitHub, Figure~\ref{fig:protection}) :

\begin{table}[H]
\centering
\small
\begin{tabularx}{\textwidth}{@{}lcL@{}}
\toprule
\textbf{Réglage} & \textbf{Valeur} & \textbf{Effet} \\
\midrule
Pull request obligatoire & oui & Aucun \emph{push} direct sur \code{main} (0 approbation requise : dépôt personnel) \\
Contrôle exigé & \textbf{Security gate} & La fusion est impossible tant que ce job n'est pas vert \\
Branche à jour avant fusion & oui & Le contrôle a bien porté sur le code qui sera fusionné \\
Inclure les administrateurs & \textbf{oui} & Personne, pas même le propriétaire, ne contourne la règle \\
\emph{Force push} et suppression & interdits & L'historique de \code{main} ne peut pas être réécrit \\
\bottomrule
\end{tabularx}
\end{table}

\figshot{t_api_protection.png}{0.82}{Règle de protection de \code{main}, telle que renvoyée par l'API GitHub.}{fig:protection}

La page des branches affiche l'icône de protection à côté de \code{main}, avec le résultat des contrôles
de chaque branche (Figure~\ref{fig:branches}).

\figshot{t_gh_branches.png}{0.92}{Page des branches : \code{main} protégée (bouclier), 8/8 contrôles réussis.}{fig:branches}

\subsection{Le job \texttt{release} : ne jamais publier une version non sécurisée}
\begin{lstlisting}[caption={Conditions de la publication}]
release:
  needs: security-gate
  if: github.event_name == 'push' && github.ref == 'refs/heads/main'
  permissions:
    contents: read
    packages: write
\end{lstlisting}

Le job télécharge l'image \textbf{testée par le DAST} et la publie sur GitHub Container Registry avec
deux étiquettes : le SHA du commit et \code{latest}. Il ne s'exécute que sur un \emph{push} sur
\code{main} et si \code{security-gate} est vert : \textbf{aucun chemin ne publie une image non
validée}. Le premier run complet sur \code{main} le démontre (Figures~\ref{fig:runmain}
et~\ref{fig:release}).

\figshot{t_gh_run_main.png}{0.92}{Premier run complet sur \code{main} : les huit jobs sont verts, dont la publication.}{fig:runmain}

\figshot{t_log_release.png}{0.95}{Journal de la publication : les deux étiquettes pointent vers le même \emph{digest}.}{fig:release}

% ---------------------------------------------------------------------------
\section{Partie V -- Validation : pull request vulnérable, puis corrigée}
% ---------------------------------------------------------------------------
\subsection{Scénario}
Un développeur ajoute une page de diagnostic réseau pour le support :
\code{/diagnostic?host=\allowbreak<nom>} affiche la résolution DNS d'un nom d'hôte. Deux pull requests
proposent cette fonctionnalité à partir de \code{main}, branche protégée (Figure~\ref{fig:pulls}). Elles ont été
ouvertes en même temps et exécutées en parallèle.

\begin{table}[H]
\centering
\small
\begin{tabularx}{\textwidth}{@{}llLl@{}}
\toprule
\textbf{Pull request} & \textbf{Branche} & \textbf{Contenu} & \textbf{Exécution} \\
\midrule
\href{https://github.com/YassineDHML/tp2-devsecops-cicd-pipeline/pull/1}{n° 1} & \code{feature/diagnostic-reseau} & Version \textbf{volontairement vulnérable} et dépendance ancienne & \#5 \\
\href{https://github.com/YassineDHML/tp2-devsecops-cicd-pipeline/pull/2}{n° 2} & \code{fix/diagnostic-securise} & Même fonctionnalité, \textbf{corrigée} & \#6 puis \#11 \\
\bottomrule
\end{tabularx}
\end{table}

\figshot{t_gh_pulls.png}{0.92}{Liste des pull requests : la n°~2 est verte (8/8), la n°~1 est en échec (3/8).}{fig:pulls}

\subsection{Pull request n°~1 : la version vulnérable est bloquée}
\begin{lstlisting}[language=Python,caption={Code de la route \code{/diagnostic} proposée par la PR n°~1}]
@app.route("/diagnostic")
def diagnostic():
    host = request.args.get("host", "localhost")
    result = subprocess.run("getent hosts " + host, shell=True, capture_output=True, text=True, timeout=10)
    return "<h1>Resolution de " + host + "</h1><pre>" + (result.stdout or "aucun resultat") + "</pre>"
\end{lstlisting}

Un lien \og Diagnostic \fg{} est ajouté dans la barre de navigation, un test vérifie le cas
nominal, et \code{requirements.txt} épingle \code{certifi==2023.5.7}. Quatre défauts sont introduits :

\begin{table}[H]
\centering
\small
\begin{tabularx}{\textwidth}{@{}rLL@{}}
\toprule
\textbf{\#} & \textbf{Défaut} & \textbf{Catégorie} \\
\midrule
1 & Route \textbf{publique} : pas de \code{@login\_required} & Contrôle d'accès défaillant (OWASP A01) \\
2 & \code{shell=True} avec \code{host} concaténé dans la commande & \textbf{Injection de commande} (CWE-78, OWASP A03) \\
3 & \code{host} recopié dans une chaîne HTML renvoyée telle quelle & \textbf{XSS réfléchi} (CWE-79, OWASP A03) \\
4 & \code{certifi==2023.5.7} & Composant vulnérable (OWASP A06) : CVE-2023-37920 (CVSS 9.8) et CVE-2024-39689 (CVSS 7.5) \\
\bottomrule
\end{tabularx}
\end{table}

\noindent\emph{Remarque :} \code{requests} 2.34.2 exige \code{certifi>=2023.5.7} ; c'est la plus ancienne
version que \code{pip} accepte, et elle est vulnérable. Un premier essai avec \code{certifi==2022.12.7}
avait échoué dès la résolution des dépendances.

\subsubsection*{Résultat du pipeline}
\begin{table}[H]
\centering
\begin{tabular}{@{}lcr@{}}
\toprule
\textbf{Job} & \textbf{Résultat} & \textbf{Durée} \\
\midrule
Build et tests unitaires & réussi (67 tests) & 25 s \\
Build de l'image Docker & réussi & 28 s \\
SAST -- Bandit & \textbf{échec} & 15 s \\
SCA -- pip-audit & \textbf{échec} & 18 s \\
SCA -- OWASP Dependency-Check & \textbf{échec} & 34 s \\
DAST -- OWASP ZAP & \textbf{échec} & 5 min 07 s \\
\textbf{Security gate} & \textbf{échec} & 4 s \\
Release & ignorée & -- \\
\bottomrule
\end{tabular}
\caption{Pipeline de la PR n°~1 (exécution \#5)}
\end{table}

Tous les contrôles se sont exécutés, puisqu'ils sont indépendants : \textbf{un seul run donne la liste
complète des problèmes}. La publication est ignorée et la protection de \code{main} bloque la fusion
(Figures~\ref{fig:pr1}, \ref{fig:runpr1} et~\ref{fig:annot}).

\figshot{t_gh_pr1.png}{0.88}{Pull request n°~1 : le dernier commit est marqué en échec, et le commentaire récapitule les contrôles.}{fig:pr1}

\figshot{t_gh_run_pr1.png}{0.92}{Exécution \#5 : build et tests réussis, SAST, SCA, DAST et porte de sécurité en échec, release ignorée.}{fig:runpr1}

\figshot{t_gh_run_pr1_annot.png}{0.92}{Annotations de l'exécution \#5 : code de sortie 1 pour Bandit, pip-audit, ZAP et la porte, code 15 pour Dependency-Check.}{fig:annot}

\subsubsection*{SAST : Bandit détecte l'injection de commande}
Le gate signale la règle \code{B602} (sévérité High, confiance High) à la ligne 287
(Figures~\ref{fig:banditlog} et~\ref{fig:banditrep}). Les deux alertes \emph{Low} sont celles de
l'application saine ; seule la High bloque.

\figshot{t_log_bandit.png}{0.95}{Journal du gate SAST : \code{B602} (High) et code de sortie 1.}{fig:banditlog}

\figshot{r_bandit_pr1.png}{0.88}{Rapport HTML de Bandit sur la PR n°~1 : \code{B602} en rouge.}{fig:banditrep}

\subsubsection*{SCA : pip-audit et Dependency-Check détectent \texttt{certifi}}
pip-audit signale 3 avis (6 lignes, chacun répété une fois par version corrigée), Dependency-Check
2 CVE de score élevé (Figures~\ref{fig:pipauditlog}, \ref{fig:dcrep} et~\ref{fig:dclog}).

\figshot{t_log_pipaudit.png}{0.9}{Journal du gate pip-audit : \code{certifi} 2023.5.7 est vulnérable.}{fig:pipauditlog}

\figshot{r_dc_pr1.png}{0.82}{Rapport Dependency-Check : \code{certifi:2023.5.7}, CVE-2023-37920 (CVSS 9.8, Critical) et CVE-2024-39689 (CVSS 7.5, High).}{fig:dcrep}

\figshot{t_log_dc.png}{0.95}{Journal du gate Dependency-Check : seuil CVSS dépassé, code de sortie 15.}{fig:dclog}

\subsubsection*{DAST : ZAP prouve l'exploitation}
ZAP relève \textbf{2 alertes High} sur \code{/diagnostic} (Figures~\ref{fig:zappr1}, \ref{fig:zapcmd}
et~\ref{fig:zapxss}) :

\figshot{r_zap_pr1.png}{0.88}{Rapport ZAP sur la PR n°~1 : 2 alertes High, en plus des 4 Medium de l'application saine.}{fig:zappr1}

\figshot{r_zap_pr1_cmdinj.png}{0.95}{\emph{Remote OS Command Injection} : ZAP a envoyé \code{localhost\&cat /etc/passwd\&} et retrouvé \code{root:x:0:0} dans la réponse.}{fig:zapcmd}

\figshot{r_zap_pr1_xss.png}{0.95}{\emph{Cross Site Scripting (Reflected)} : la charge utile est recopiée telle quelle dans la page.}{fig:zapxss}

\figshot{t_log_dast_gate.png}{0.95}{Journal du gate DAST sur la PR n°~1 : deux alertes High, code de sortie 1.}{fig:dastgate}

\noindent Le contenu de \code{/etc/passwd} dans la réponse est la \textbf{preuve} que la commande a été
exécutée par le serveur. La même attaque, rejouée à la main sur l'image de la PR n°~1, donne
la Figure~\ref{fig:exploit} ; le HTML injecté est interprété par le navigateur (Figure~\ref{fig:xss}). La
politique CSP de l'application empêche l'exécution d'un script \emph{inline}, mais pas l'injection de
balises.

\figshot{a_vuln_exploit.png}{0.8}{Exploitation manuelle de la PR n°~1 : le serveur renvoie \code{/etc/passwd}.}{fig:exploit}

\figshot{a_vuln_xss.png}{0.8}{Injection de HTML dans la réponse : le titre rouge n'existe pas dans le code de la page.}{fig:xss}

\subsubsection*{La porte de sécurité bloque la fusion}
\begin{lstlisting}[basicstyle=\ttfamily\footnotesize,caption={Extrait de la sortie de la porte de sécurité sur la PR n°~1}]
| Controle | Resultat |
| build-test | success |
| build-image | success |
| sast | failure |
| sca-pip-audit | failure |
| sca-dependency-check | failure |
| dast | failure |
##[error]Au moins un controle n'est pas vert : version non securisee, publication bloquee.
\end{lstlisting}

\figshot{t_log_gate.png}{0.88}{Journal de la porte de sécurité sur la PR n°~1.}{fig:gate}

\figshot{t_api_merge.png}{0.82}{État de fusion renvoyé par l'API : \code{BLOCKED} pour la PR n°~1, \code{CLEAN} pour la PR n°~2.}{fig:merge}

\subsubsection*{Ce que chaque contrôle a vu}
\begin{table}[H]
\centering
\small
\begin{tabularx}{\textwidth}{@{}LccccX@{}}
\toprule
\textbf{Défaut de la PR n°~1} & \textbf{Tests} & \textbf{Bandit} & \textbf{pip-audit} & \textbf{Dep.-Check} & \textbf{ZAP} \\
\midrule
1. Route publique & non vu & non vu & -- & -- & non signalé en tant que tel \\
2. Injection de commande & non vue & \textbf{détectée} (B602) & -- & -- & \textbf{détectée} (High) \\
3. XSS réfléchi & non vu & \textbf{non vu} & -- & -- & \textbf{détecté} (High) \\
4. \code{certifi} vulnérable & non vue & -- & \textbf{détectée} & \textbf{détectée} & -- \\
\bottomrule
\end{tabularx}
\caption{Matrice défauts $\times$ contrôles}
\end{table}

\begin{enumerate}[leftmargin=*]
  \item \textbf{Les tests unitaires sont verts.} Ils vérifient que la fonctionnalité marche, pas qu'elle
        est sûre : seuls les contrôles de sécurité arrêtent cette PR.
  \item \textbf{SAST et DAST se complètent.} L'injection de commande est vue par les deux : Bandit
        repère le motif \code{shell=True}, ZAP prouve l'exécution. Le XSS n'est vu \textbf{que par le
        DAST} : \code{"<h1>" + host + "</h1>"} ne contient aucun motif que Bandit sache reconnaître.
  \item \textbf{Les deux outils SCA concordent} sur le même composant, avec des vocabulaires
        différents (avis PyPA d'un côté, CVE et score CVSS de l'autre).
  \item \textbf{Aucun outil ne signale l'absence d'authentification en tant que telle} : le contrôle
        d'accès relève de la logique métier. Cette règle doit être verrouillée par un \textbf{test
        unitaire}, ce que fait la PR n°~2.
\end{enumerate}

\subsection{Pull request n°~2 : la version corrigée passe}
\begin{lstlisting}[language=Python,caption={Code corrigé de la route \code{/diagnostic} (PR n°~2)}]
HOSTNAME_RE = re.compile(
    r"(?=.{1,253}\Z)(?!-)[A-Za-z0-9-]{1,63}(?<!-)(\.(?!-)[A-Za-z0-9-]{1,63}(?<!-))*"
)

@app.route("/diagnostic")
@login_required
def diagnostic():
    host = request.args.get("host", "localhost")
    if not HOSTNAME_RE.fullmatch(host):
        abort(400, "Nom d'hote invalide")
    try:
        infos = socket.getaddrinfo(host, None, proto=socket.IPPROTO_TCP)
    except socket.gaierror:
        output = "Aucune adresse trouvee"
    else:
        output = "\n".join(sorted({info[4][0] for info in infos}))
    return render_template("tool.html", title="Resolution de " + host, output=output)
\end{lstlisting}

\begin{table}[H]
\centering
\small
\begin{tabularx}{\textwidth}{@{}LL@{}}
\toprule
\textbf{Défaut de la PR n°~1} & \textbf{Correction} \\
\midrule
Route publique & \code{@login\_required} : redirection vers \code{/login} sans session ; lien affiché aux seuls utilisateurs connectés \\
Injection de commande & \textbf{Plus aucun sous-processus} : résolution avec \code{socket.getaddrinfo}. Le danger est supprimé à la source plutôt que filtré \\
Entrée non contrôlée & Nom d'hôte validé par \textbf{liste blanche} (RFC 1123, \code{fullmatch}), erreur 400 sinon \\
XSS réfléchi & Réponse rendue par un \textbf{template Jinja2} (auto-échappement) au lieu d'une chaîne HTML concaténée \\
\code{certifi} ancien & Pas d'épinglage : \code{main} garde \code{certifi==2026.10.23} \\
\bottomrule
\end{tabularx}
\end{table}

Les tests ajoutés verrouillent ces propriétés : authentification obligatoire, 9 noms d'hôte dangereux
rejetés (\code{;}, \code{\$( )}, \code{\&\&}, \code{|}, \code{-c1}, balise \code{<script>}, saut de ligne,
300 caractères), aucun appel à \code{subprocess}, hôte inconnu.

\subsubsection*{Résultat du pipeline}
\begin{table}[H]
\centering
\begin{tabular}{@{}lcr@{}}
\toprule
\textbf{Job} & \textbf{Résultat} & \textbf{Durée} \\
\midrule
Build et tests unitaires & réussi (78 tests) & 26 s \\
Build de l'image Docker & réussi & 28 s \\
SAST -- Bandit & réussi (0 High, 2 Low acceptées) & 17 s \\
SCA -- pip-audit & réussi (0 vulnérabilité) & 17 s \\
SCA -- OWASP Dependency-Check & réussi (0 CVE) & 38 s \\
DAST -- OWASP ZAP & réussi (\textbf{0 High}, 4 Medium) & 4 min 57 s \\
\textbf{Security gate} & \textbf{réussi} & 3 s \\
Release & ignorée sur une PR & -- \\
\bottomrule
\end{tabular}
\caption{Pipeline de la PR n°~2 (exécution \#6, commit \code{c93889e})}
\end{table}

L'état de la PR est \code{CLEAN} (Figure~\ref{fig:merge}) : toutes les règles de la protection de
\code{main} sont satisfaites et la fusion est autorisée. L'historique de la PR (Figure~\ref{fig:pr2})
montre aussi l'épisode du faux positif de la \S\ref{sec:fp} : les croix intermédiaires sont des
exécutions annulées par une poussée plus récente, puis le run bloqué, avant le retour au vert.

\figshot{t_gh_pr2.png}{0.88}{Pull request n°~2 : 6 commits, premier et dernier commits validés.}{fig:pr2}

\figshot{t_gh_run_pr2.png}{0.92}{Exécution \#11 de la PR n°~2 : tous les contrôles sont verts, release ignorée (événement \code{pull\_request}).}{fig:runpr2}

\noindent Les trois scanners n'ont plus rien à signaler (Figures~\ref{fig:banditpr2} et~\ref{fig:dcpr2}, et
Figure~\ref{fig:zappr2} pour ZAP) :

\figshot{r_bandit_pr2.png}{0.88}{Rapport Bandit sur la PR n°~2 : plus de \code{B602}, seules restent les deux alertes Low acceptées.}{fig:banditpr2}

\figshot{r_dc_pr2.png}{0.82}{Rapport Dependency-Check sur la PR n°~2 : 14 dépendances analysées, aucune vulnérable.}{fig:dcpr2}

\subsubsection*{Comportement de l'application corrigée}
Sans session, la route redirige vers la page de connexion (Figure~\ref{fig:secure1}). Connecté, un nom
d'hôte invalide est refusé (HTTP 400, Figure~\ref{fig:secure2}) et un nom valide renvoie ses adresses
(Figure~\ref{fig:secure3}).

\figshot{a_secure_login_redirect.png}{0.75}{Version corrigée : accès sans session, redirection vers la page de connexion.}{fig:secure1}

\figshot{a_secure_400.png}{0.75}{Version corrigée : la tentative \code{localhost; id} est refusée (HTTP 400).}{fig:secure2}

\figshot{a_secure_diag_ok.png}{0.75}{Version corrigée : un nom d'hôte valide renvoie ses adresses.}{fig:secure3}

\subsection{Publication de l'image}
Le job \code{release} est \textbf{ignoré sur une pull request} (\code{github.event\_name == 'push'} est
faux) : c'est ce que l'on observe sur les deux PR. Il s'exécutera après la fusion, sur le \emph{push}
résultant sur \code{main}, et seulement si la porte de sécurité est verte, comme l'a démontré le run de
la Figure~\ref{fig:runmain}.

\subsection{Synthèse comparée}
\begin{table}[H]
\centering
\small
\begin{tabularx}{\textwidth}{@{}lLL@{}}
\toprule
\textbf{Contrôle} & \textbf{PR n°~1 (vulnérable)} & \textbf{PR n°~2 (corrigée)} \\
\midrule
Tests unitaires & 67 réussis & 78 réussis \\
Bandit & 1 High, 2 Low : \textbf{bloque} & 0 High, 2 Low : passe \\
pip-audit & 3 avis sur \code{certifi} : \textbf{bloque} & 0 : passe \\
Dependency-Check & 1 Critical (9.8), 1 High (7.5) : \textbf{bloque} & 0 CVE : passe \\
ZAP & 2 High, 4 Medium : \textbf{bloque} & 0 High, 4 Medium : passe \\
Porte de sécurité & \textbf{échec}, état \code{BLOCKED} & réussie, état \code{CLEAN} \\
Publication & ignorée & ignorée (PR), exécutée après la fusion \\
\bottomrule
\end{tabularx}
\caption{Bilan comparé des deux pull requests}
\end{table}

% ---------------------------------------------------------------------------
\section{Conclusion}
% ---------------------------------------------------------------------------
Ce TP a transformé les contrôles ponctuels du TP1 en un \textbf{garde-fou permanent}. Chaque
modification traverse un pipeline qui compile, teste, analyse le code (Bandit), les dépendances
(pip-audit, Dependency-Check) et l'application en cours d'exécution (OWASP ZAP sur un environnement
éphémère). Un job unique, \code{Security gate}, agrège les résultats ; la branche \code{main} est
protégée, y compris pour l'administrateur, et l'image n'est publiée que si tout est vert. La
validation l'a démontré : la pull request vulnérable, dont les tests unitaires passaient, a été arrêtée
par quatre contrôles distincts, alors que la version corrigée est passée sans alerte bloquante.

Le TP met en évidence plusieurs points :
\begin{itemize}
  \item \textbf{Des tests fonctionnels verts ne disent rien de la sécurité.}
  \item \textbf{SAST et DAST se complètent} : le XSS n'est vu que par le DAST ; l'injection de commande
        l'est par les deux.
  \item \textbf{La SCA doit être continue} : \code{certifi} est devenu vulnérable sans qu'une ligne de
        code change, et deux outils SCA ne voient pas les mêmes choses.
  \item \textbf{Un garde-fou doit échouer fermé} : rapport manquant, contrôle ignoré, scan vide, tous
        ces cas font échouer le pipeline. Un contrôle obligatoire \emph{ignoré} passe pour réussi dans
        la protection de branche, d'où le \code{if: always()} de la porte de sécurité.
  \item \textbf{Un DAST produit des faux positifs} : il faut un processus de triage, pas un
        \og relancer jusqu'à ce que ça passe \fg{} ni une baisse du seuil.
\end{itemize}

\paragraph{Limites.} La couverture du DAST dépend de ce que ZAP connaît (les routes non liées doivent
être déclarées) ; aucun scanner n'est exhaustif (Dependency-Check n'a pas vu l'avis \code{certifi}, Bandit
ne voit pas le XSS) ; le pipeline lui-même est une cible : une PR s'exécute avec sa propre version de
\code{ci.yml}, d'où l'intérêt d'un \code{CODEOWNERS} imposant une revue sur \code{.github/}.

\paragraph{Pistes d'amélioration.} Corriger les alertes Medium du DAST (protection CSRF avec
Flask-WTF, CSP complète, HTTPS), Dependabot pour les dépendances et les actions, analyse de l'image avec
Trivy, recherche de secrets avec gitleaks, SBOM et signature de l'image (cosign), approbation manuelle
avant la publication, scan DAST complet nocturne.

% ---------------------------------------------------------------------------
\bigskip
\section*{Annexe A -- Fichier \texttt{.github/workflows/ci.yml}}\label{sec:ci}
\addcontentsline{toc}{section}{Annexe A -- Fichier ci.yml}
% ---------------------------------------------------------------------------
\lstinputlisting[basicstyle=\ttfamily\scriptsize,caption={Pipeline complet (fichier \code{.github/workflows/ci.yml})}]{../../.github/workflows/ci.yml}

% ---------------------------------------------------------------------------
\section*{Annexe B -- Plan ZAP \texttt{zap/automation.yaml}}\label{sec:zapplan}
\addcontentsline{toc}{section}{Annexe B -- Plan ZAP}
% ---------------------------------------------------------------------------
\lstinputlisting[basicstyle=\ttfamily\scriptsize,caption={Plan d'automatisation ZAP}]{../../zap/automation.yaml}

% ---------------------------------------------------------------------------
\section*{Annexe C -- Index des preuves d'exécution}
\addcontentsline{toc}{section}{Annexe C -- Index des preuves}
% ---------------------------------------------------------------------------
Les rapports et journaux des exécutions de référence sont conservés dans le dépôt, dossier
\code{reports/} (les artefacts GitHub expirent après 90 jours).

\begin{table}[H]
\centering
\small
\begin{tabularx}{\textwidth}{@{}llL@{}}
\toprule
\textbf{Dossier} & \textbf{Exécution} & \textbf{Contenu} \\
\midrule
\code{pr1-vulnerable/} & \#5, PR n°~1 (\code{0285a71}) & Pipeline en échec : Bandit, pip-audit, Dependency-Check, ZAP, journaux des étapes en échec \\
\code{pr2-fixed/} & \#6, PR n°~2 (\code{c93889e}) & Pipeline vert : mêmes rapports, sans alerte bloquante \\
\code{pr2-faux-positif/} & \#10, PR n°~2 (\code{fe474eb}) & DAST bloqué par le faux positif ZAP \\
\code{main-release/} & \#3, \emph{push} sur \code{main} (\code{2b88fac}) & Résultat des huit jobs et journal de la publication \\
\bottomrule
\end{tabularx}
\end{table}

\noindent Les figures de ce compte rendu proviennent de ces exécutions : pages GitHub du dépôt (\code{t\_gh\_*}),
rapports HTML des outils (\code{r\_*}), journaux des étapes (\code{t\_log\_*}) et application exécutée en
local (\code{a\_*}).

\end{document}
